%-----------------------------------------------------------------------------%
\chapter{\babTiga}
\label{cha:metodologi}
%-----------------------------------------------------------------------------%
Bab ini menjelaskan metodologi penelitian dan rancangan sistem DroneVL. Rancangan memisahkan inferensi VLM atau VLA yang bergantung pada \textit{backend} dari komponen navigasi yang dibekukan pada seluruh kondisi eksperimen. Spesifikasi berikut menjadi acuan implementasi, pencatatan, dan evaluasi; setiap nilai konfigurasi yang belum bersifat universal dibekukan dalam manifest eksperimen sebelum pengujian dimulai.

%-----------------------------------------------------------------------------%
\section{Desain Penelitian}
\label{sec:desain_penelitian}
%-----------------------------------------------------------------------------%
Penelitian menggunakan pendekatan rekayasa sistem dan evaluasi eksperimental. Artefak yang dibangun meliputi DroneVL sebagai adaptor untuk VLM dan VLA, serta DroneVL Arena sebagai kerangka evaluasi \textit{closed-loop}. Evaluasi membedakan dua kondisi. Kondisi antarmuka bersama menyetarakan episode, sensor yang diizinkan, resolusi masukan, riwayat, \textit{prompt}, frekuensi keputusan, batas waktu, kebijakan coba-ulang, pengendali, dan evaluator sejauh didukung semua \textit{backend}. Kondisi kapabilitas asli mencatat konfigurasi terbaik yang sah bagi setiap \textit{backend}; hasilnya dilaporkan terpisah dan tidak ditafsirkan sebagai pengaruh model semata.

Eksperimen dibagi menjadi tiga keluarga. Eksperimen VLM membandingkan GPT-5.6, Qwen3.6 35B-A3B, dan Gemini sebagai perencana semantik. Eksperimen portabilitas VLA membandingkan CognitiveDrone, OpenVLA-7B yang diadaptasi untuk UAV, dan OpenPI $\pi_{0.5}$ yang diadaptasi untuk UAV. Eksperimen utama adaptasi membandingkan model dasar OpenVLA-7B (M0), OpenVLA-7B dengan LoRA yang diadaptasi untuk UAV (M1), dan M1 yang dilanjutkan GRPO (M2). Hasil ketiga keluarga dilaporkan terpisah; perbandingan lintas keluarga hanya bersifat deskriptif karena representasi aksi dan riwayat pelatihan berbeda.

Setiap episode memiliki pasangan percobaan yang sama di dalam keluarga eksperimen yang relevan. Kesamaan jalur hilir mengendalikan perubahan pengendali dan skenario, tetapi tidak menghapus perbedaan bawaan seperti modalitas, panjang konteks, atau latensi layanan. Karena itu, konfigurasi efektif dan kehilangan informasi akibat normalisasi dicatat bersama hasil eksperimen.

%-----------------------------------------------------------------------------%
\section{Rancangan Arsitektur DroneVL}
\label{sec:rancangan_dronevl}
%-----------------------------------------------------------------------------%
Arsitektur terdiri atas komponen pembentuk pengamatan, adaptor \textit{backend}, pengurai respons, pengelola misi, gerbang keselamatan, pengendali tingkat rendah, dan pencatat eksperimen. Alur representasi dibekukan sebagai berikut: \textit{CanonicalObservation} $\rightarrow$ permintaan khusus-\textit{backend} $\rightarrow$ respons mentah dan metadata $\rightarrow$ hasil deserialisasi khusus-\textit{backend} $\rightarrow$ aksi kanonis atau usulan peristiwa misi $\rightarrow$ keputusan keselamatan $\rightarrow$ aksi aktual. Adaptor hanya bertanggung jawab atas dua tahap khusus-\textit{backend}; pengurai memiliki satu tanggung jawab untuk membentuk representasi kanonis dan melaporkan kegagalan penguraian.

Pada jalur eksperimen VLM, GPT-5.6, Qwen3.6 35B-A3B, dan Gemini hanya terhubung ke jalur penerbangan melalui DroneVL Model Adapter. GPT-5.6 menggunakan antarmuka API gambar--teks, Qwen3.6 35B-A3B menggunakan antarmuka gambar--teks dengan arsitektur \textit{mixture-of-experts} (MoE) multimodal, sedangkan Gemini menggunakan Live API dua arah untuk menerima konteks multimodal dan mengembalikan teks atau pemanggilan fungsi. Respons \textit{streaming} tersebut tidak diteruskan langsung ke UAV; respons harus dihimpun pada batas keputusan, dideserialisasi, lalu dinormalisasi menjadi aksi tingkat tinggi yang dibatasi atau ditolak sebelum memasuki gerbang keselamatan.

Jalur VLA menggunakan kontrak pengamatan kanonis yang sama, tetapi adaptor VLA juga membaca deskriptor perwujudan. Deskriptor tersebut memuat skema aksi, satuan, kerangka koordinat, horizon, dan versi \textit{checkpoint}. CognitiveDrone, OpenVLA-7B yang diadaptasi untuk UAV, dan OpenPI $\pi_{0.5}$ yang diadaptasi untuk UAV mengembalikan \textit{NativeActionEnvelope} berversi. Adaptor mendeserialisasi amplop tersebut ke kandidat aksi kanonis, tanpa meneruskan aksi asli langsung ke UAV.

OpenVLA-7B menggunakan token aksi UAV yang dipelajari dari split pelatihan. OpenPI $\pi_{0.5}$ menggunakan potongan aksi UAV kontinu yang didekode per horizon, sedangkan CognitiveDrone menggunakan skema aksi empat-dimensi yang ditetapkan dalam manifest. Semua kandidat selanjutnya melewati pengurai dan gerbang keselamatan yang sama.

Perbedaan bentuk masukan dan keluaran ketiga model sebelum normalisasi menjadi objek pencatatan pada lapisan adaptor, bukan alasan untuk mengubah pengendali, evaluator, atau aturan terminal.

\subsection{Kontrak Pengamatan \textit{CanonicalObservation}}
\label{subsec:canonical_observation}
Setiap keputusan menerima \texttt{CanonicalObservation} berversi yang sedikitnya memuat \texttt{schema\_version}, \texttt{episode\_id}, \texttt{decision\_id}, waktu simulator, waktu dinding, dan cap waktu penangkapan sensor. Isinya terdiri atas citra RGB depan; kedalaman terdaftar beserta satuan dan kerangka koordinat kamera; keadaan UAV yang memuat pose, kecepatan, dan orientasi; instruksi; serta ringkasan kemajuan misi yang dihasilkan pengelola misi. Manifest menentukan resolusi, pemotongan, normalisasi, usia maksimum frame, frekuensi keputusan, dan cara riwayat dipadatkan.

Segmentasi instans, label objek, transformasi objek, serta anotasi target tidak berada dalam \texttt{CanonicalObservation}. Data tersebut hanya dapat dibaca pembangun manifest, \textit{baseline} oracle, dan evaluator melalui kanal oracle terpisah. Pembangun pengamatan mengunci pasangan RGB, kedalaman, dan odometri pada waktu simulator yang sama; frame yang melampaui usia maksimum dibuang, bukan dikirimkan ke \textit{backend}.

\subsection{Kontrak Aksi dan Peristiwa Misi}
\label{subsec:canonical_action}
Aksi kendali kanonis memuat \texttt{action\_type}, \texttt{observation\_id}, \texttt{sequence\_id}, \texttt{valid\_until}, dan parameter bernilai hingga. \texttt{MOVE\_RELATIVE} memuat perpindahan $(\Delta x,\Delta y,\Delta z)$ dalam meter pada kerangka badan UAV yang dinyatakan manifest; \texttt{ROTATE\_YAW} memuat perubahan yaw dalam radian; \texttt{HOVER\_AND\_OBSERVE} memuat durasi terbatas; dan \texttt{ABORT} tidak memiliki parameter gerak. Satu konversi kerangka dari aksi kanonis ke perintah pengendali dilakukan pada batas pengendali, sehingga adaptor tidak dapat melakukan transformasi koordinat tambahan.

\texttt{DECLARE\_TARGET} dan \texttt{ADVANCE\_SUBGOAL} bukan aksi kendali. Keduanya adalah usulan peristiwa misi yang memuat identitas atau predikat yang diklaim; pengelola misi memverifikasi usulan terhadap kondisi episode dan kanal oracle evaluator sebelum mengubah keadaan misi. Penolakan usulan dicatat sebagai kegagalan semantik, bukan sebagai kegagalan pengendali.

Pengurai memeriksa skema, bidang wajib, tipe data, nilai numerik berhingga, rentang aksi, cap waktu, dan legalitas aksi pada keadaan misi saat ini. Satu \textit{prompt} perbaikan dapat digunakan untuk respons tidak valid. Setelah $N_{\mathrm{parse}}$ kegagalan penguraian berturut-turut, dengan nilai yang dibekukan dalam manifest, episode diakhiri dengan \texttt{PARSE\_FAILURE}; penghitung hanya direset oleh aksi kanonis yang valid.

Untuk VLA, \textit{NativeActionEnvelope} wajib memuat \texttt{action\_schema\_version}, \texttt{checkpoint\_id}, \texttt{native\_action}, \texttt{coordinate\_frame}, \texttt{unit}, dan \texttt{horizon}. Pemetaan ke \texttt{MOVE\_RELATIVE}, \texttt{ROTATE\_YAW}, \texttt{HOVER\_AND\_OBSERVE}, atau \texttt{ABORT} ditetapkan per \textit{backend} dalam manifest dan diuji pada split pengembangan. Tidak ada pemetaan implisit dari aksi sendi robot, gripper, atau ruang aksi manipulasi ke aksi UAV. Apabila skema, \textit{checkpoint}, atau pemetaan tidak cocok, respons ditolak sebagai \texttt{VLA\_SCHEMA\_FAILURE} dan tidak dieksekusi.

%-----------------------------------------------------------------------------%
\section{Rancangan Lingkungan DroneVL Arena}
\label{sec:rancangan_arena}
%-----------------------------------------------------------------------------%
Setiap episode DroneVL Arena didefinisikan oleh manifest berversi yang memuat identitas dan versi adegan, \textit{seed}, pose awal, instruksi, target, kondisi lingkungan, anggaran waktu dan keputusan, batasan keselamatan, kondisi terminal, serta anotasi \textit{ground truth}. Setiap episode menerima satu atau lebih label kemampuan: penambatan visual, penalaran spasial, pencarian objek, navigasi semantik, dan instruksi multilangkah. Hasil dilaporkan menurut hasil misi akhir dan menurut label kemampuan; satu episode multikemampuan tidak dipaksa masuk ke satu kategori eksklusif.

Manifest membedakan split pelatihan, pengembangan, dan pengujian. Split pengujian dibekukan sebelum penyetelan dan dipisahkan dari pelatihan maupun pengembangan berdasarkan tata ruang, instans objek, komposisi instruksi, dan kondisi visual. Keberhasilan misi memerlukan pemenuhan target, predikat relasi bila ada, urutan subtujuan, toleransi pose, dan durasi bertahan yang ditetapkan oleh manifest. Untuk relasi spasial, evaluator memeriksa kecocokan instans target dan predikat pose secara terpisah dengan kerangka acuan, toleransi, serta waktu penilaian yang eksplisit.

%-----------------------------------------------------------------------------%
\section{Protokol Eksperimen, Log, dan Metrik}
\label{sec:protokol_eksperimen}
%-----------------------------------------------------------------------------%
Sebelum pengujian VLM, konfigurasi GPT-5.6, Qwen3.6 35B-A3B dengan tag \texttt{qwen3.6:35b-a3b}, dan Gemini dikunci serta dicatat. Untuk Gemini, catatan juga memuat mode sesi Live API, modalitas \textit{streaming}, batas waktu pengumpulan respons pada setiap keputusan, serta aturan pemetaan pemanggilan fungsi atau teks ke respons mentah.

Sebelum pengujian portabilitas VLA, konfigurasi CognitiveDrone, OpenVLA-7B yang diadaptasi untuk UAV, dan OpenPI $\pi_{0.5}$ yang diadaptasi untuk UAV dikunci serta dicatat. Catatan tersebut memuat asal \textit{checkpoint}, split dan versi trajektori adaptasi, skema atau tokenisasi aksi UAV, horizon keluaran, serta prosedur dekode. Untuk seluruh \textit{backend}, catatan juga memuat identitas versi atau rilis, penyedia dan jenis antarmuka, kapabilitas yang dideklarasikan, prapemrosesan gambar, resolusi efektif, \textit{prompt}, parameter dekode, \textit{runtime}, perangkat keras lokal, dan wilayah atau konfigurasi layanan apabila relevan. Seluruh \textit{backend} pada kondisi antarmuka bersama menggunakan episode, pengendali, batasan keselamatan, kebijakan coba-ulang, dan evaluator yang sama.

Log per keputusan memuat identitas episode dan manifest, versi anotasi, \texttt{observation\_id}, respons mentah, hasil deserialisasi, aksi kanonis yang diusulkan, usulan peristiwa misi, keputusan gerbang, aksi aktual, keadaan UAV, usia pengamatan, waktu inferensi, waktu keputusan ujung-ke-ujung, dan alasan terminal. Untuk penambatan, log juga menyimpan target prediksi, target acuan, aturan pencocokan, hasil kecocokan instans, dan hasil predikat relasi. Dengan demikian, metrik agregat dapat dihitung ulang dari log mentah.

Akurasi penambatan instans didefinisikan sebagai $A_{\mathrm{inst}}=N^{-1}\sum_{i=1}^{N}\mathbf{1}[p_i=g_i]$ pada semua peluang deklarasi target yang dievaluasi. Akurasi relasi dihitung terpisah pada peluang yang memuat predikat spasial. Tingkat keberhasilan misi adalah proporsi episode yang memenuhi seluruh kondisi terminal berhasil. Tingkat tabrakan adalah proporsi episode yang mengalami sedikitnya satu tabrakan; jumlah tabrakan per keputusan dilaporkan sebagai metrik tambahan. Tingkat aksi tidak valid adalah proporsi respons awal yang gagal diurai, sedangkan kegagalan setelah \textit{prompt} perbaikan dicatat terpisah. Tingkat intervensi keselamatan adalah proporsi aksi kanonis yang dipotong, diganti, atau ditolak. Efisiensi lintasan adalah rasio panjang lintasan referensi yang layak terhadap panjang lintasan aktual untuk episode berhasil, dengan lintasan referensi dibekukan oleh manifest. Latensi inferensi dan latensi ujung-ke-ujung dilaporkan sebagai median, persentil ke-95, dan distribusi per episode.

%-----------------------------------------------------------------------------%
\section{Eksperimen Adaptasi LoRA dan GRPO}
\label{sec:lora_grpo_experiment}
%-----------------------------------------------------------------------------%
LoRA dan GRPO merupakan tujuan utama penelitian pada OpenVLA-7B yang diadaptasi untuk UAV. Tiga kondisi dibandingkan: model dasar OpenVLA-7B dengan tokenisasi aksi UAV yang ditetapkan (M0); M0 dengan LoRA terawasi menggunakan trajektori CoSyS-AirSim pada split pelatihan (M1); dan M1 yang dilanjutkan GRPO menggunakan \textit{rollout} CoSyS-AirSim (M2). Data pelatihan hanya berasal dari split pelatihan; split pengembangan digunakan untuk memilih hiperparameter; split pengujian yang telah dibekukan tidak digunakan untuk penyetelan maupun \textit{rollout}.

LoRA dan GRPO dibatasi pada VLA OpenVLA-7B, bukan diterapkan pada seluruh model dalam penelitian. Tujuannya adalah mengisolasi pengaruh adaptasi kebijakan aksi UAV melalui perbandingan M0--M1--M2 dengan data, pengendali, episode, dan metrik yang sama. OpenVLA-7B dipilih karena keluarannya berupa token aksi diskret dan prosedur LoRA tersedia untuk tugas maupun perwujudan robot baru \citep{kim2024openvla}; kondisi ini memungkinkan log-probabilitas token aksi dicatat pada setiap \textit{rollout} untuk perhitungan GRPO. Kelompok VLM tetap dievaluasi sebagai perencana semantik melalui kontrak pengamatan dan aksi yang sama. Penerapan LoRA atau GRPO pada VLM akan menambahkan pertanyaan penelitian lain tentang adaptasi perencana, sekaligus mengaburkan apakah perubahan hasil berasal dari model perencana, kebijakan aksi, atau lapisan adaptor.

CognitiveDrone dan OpenPI $\pi_{0.5}$ tetap digunakan pada eksperimen portabilitas VLA, tetapi bukan sebagai kondisi LoRA+GRPO. CognitiveDrone menjadi pembanding kebijakan khusus UAV yang sudah tersedia. OpenPI $\pi_{0.5}$ dapat diadaptasi terhadap trajektori UAV sesuai prosedur pelatihannya sendiri, tetapi keluaran \textit{flow-matching}-nya berupa potongan aksi kontinu \citep{physicalintelligence2025openpi}. Karena representasi tersebut tidak serta-merta menyediakan rasio probabilitas kebijakan yang diperlukan GRPO, memasukkannya ke M0--M1--M2 akan menghasilkan perbandingan metode optimisasi yang tidak sepadan. Dengan pembatasan ini, eksperimen portabilitas menjawab kesesuaian antarmuka VLA, sedangkan eksperimen OpenVLA-7B menjawab dampak LoRA dan GRPO secara terisolasi.

Manifest pelatihan menetapkan data dan versi \textit{prompt}, tokenisasi serta skema aksi UAV, lapisan dan peringkat LoRA, kuantisasi, perangkat keras, anggaran token, jumlah \textit{rollout}, ukuran kelompok GRPO, kebijakan acuan, dan \textit{seed}. Pada M2, sejumlah \textit{rollout} bertitik awal sama disampel dari token aksi OpenVLA, log-probabilitas kebijakan lama disimpan, dan penghargaan relatif kelompok dihitung dari keberhasilan misi, kemajuan subtujuan, keselamatan, kepatuhan format aksi, serta penalti aksi tidak valid. Bobot komponen serta setiap ablasi dibekukan sebelum pelatihan. Ketiga kondisi dievaluasi dengan manifest episode, pengendali, batas keselamatan, dan metrik yang sama, sehingga perubahan hasil dapat ditafsirkan sebagai efek LoRA dan GRPO pada OpenVLA-7B UAV yang diuji. OpenPI $\pi_{0.5}$ tidak dimasukkan ke klaim GRPO; ia hanya dibandingkan pada eksperimen portabilitas VLA setelah adaptasi UAVnya sendiri.
